\documentclass[a4paper,11pt]{article}
\usepackage{exam-style-341}
\examinehead{341 农业综合知识三 · 网络技术分章练习}{第04章\quad 数据链路层}
\begin{document}

\unitcover{341 农业综合知识三 · 网络技术分章练习}{第04章\quad 数据链路层}{练习（共 100 分）}{约 60 分钟}{链路与数据链路；数据链路层的三个基本问题（封装成帧、透明传输、差错检测）；字节填充与零比特填充；CRC 循环冗余检验；停止等待协议、连续 ARQ（回退 N 帧）、选择重传 ARQ 与滑动窗口；PPP 协议的组成、帧格式与工作状态。}

\begin{knowmap}
\textbf{两个概念：}链路是一条无源的点到点物理线路段；数据链路则是链路加上实现通信协议的硬件和软件\\[2pt]
\textbf{三个基本问题：}封装成帧（帧定界）· 透明传输（字节填充、零比特填充）· 差错检测（CRC 循环冗余检验）\\[2pt]
\textbf{透明传输：}异步传输用字节填充（转义符 0x7D）；同步传输用零比特填充（每 5 个连续 1 后填 1 个 0）\\[2pt]
\textbf{CRC：}收发双方约定生成多项式 $G(x)$，发送端用模 2 除法求余数作为帧检验序列 FCS 附在帧尾，接收端用同一 $G(x)$ 检验\\[2pt]
\textbf{可靠传输：}停止等待协议（发一帧等一个确认，超时重传，1 位序号）$\rightarrow$ 连续 ARQ/回退 N 帧（发送窗口 $>1$、接收窗口 $=1$）$\rightarrow$ 选择重传（收发窗口均 $>1$，只重传出错帧）\\[2pt]
\textbf{滑动窗口：}用发送窗口和接收窗口控制可连续发送的帧数；序号为 $n$ 位时，回退 N 帧发送窗口 $\leq 2^n-1$，选择重传窗口 $\leq 2^{n-1}$\\[2pt]
\textbf{PPP：}点对点协议，由封装成帧、链路控制协议 LCP、网络控制协议 NCP 三部分组成；帧以 0x7E 定界
\end{knowmap}

\begin{keyterms}
\textbf{链路}：从一个节点到相邻节点的一段物理线路（有线或无线），中间没有任何其他交换节点。\\
\textbf{数据链路}：在链路之上加上实现通信协议的硬件和软件，也称逻辑链路。\\
\textbf{封装成帧}：在一段数据前后分别添加首部和尾部，构成一个帧，以便接收端识别帧的边界。\\
\textbf{透明传输}：无论数据部分出现什么比特组合，都能按原样正确传输，不会与帧定界符混淆。\\
\textbf{差错检测}：通过校验码发现传输过程中出现的比特差错，常用循环冗余检验 CRC。\\
\textbf{CRC（循环冗余检验）}：用生成多项式作模 2 除法，把余数作为帧检验序列的差错检测方法。\\
\textbf{FCS（帧检验序列）}：发送端由 CRC 计算出的、附加在帧尾的检验序列，供接收端检验。\\
\textbf{停止等待协议}：每发送一个帧就停止发送，等待对方确认，收到确认后再发下一帧；超时未收到确认则重传。\\
\textbf{连续 ARQ（回退 N 帧）}：发送方可连续发送多个帧，接收窗口为 1；某帧出错时重传该帧及其后所有帧。\\
\textbf{选择重传 ARQ}：接收窗口大于 1，只重传真正出错的帧，其他正确收到的帧先缓存。\\
\textbf{滑动窗口}：用窗口大小限制已发送但未确认的帧数，实现流量控制与可靠传输。\\
\textbf{PPP（点对点协议）}：两个相邻节点之间传送数据的数据链路层协议，由封装成帧、LCP、NCP 三部分组成。\\
\textbf{LCP（链路控制协议）}：PPP 中用于建立、配置、测试和释放数据链路连接的协议。\\
\textbf{NCP（网络控制协议）}：PPP 中用于配置网络层协议（如为接入主机分配 IP 地址）的协议。
\end{keyterms}

\begin{contrast}{停止等待 / 连续 ARQ（回退 N）/ 选择重传（每行依次为：发送窗口·接收窗口·重传策略）}
\citem{停止等待}{$1$}{$1$}{超时未收到确认即重传该帧，序号用 1 位即可}
\citem{连续 ARQ（回退 N）}{$>1$，且 $\leq 2^n-1$}{$1$}{某帧出错后，重传该帧及其后所有已发帧}
\citem{选择重传}{$>1$}{$>1$，且 $\leq 2^{n-1}$}{只重传出错的帧，其余正确帧缓存后按序交付}
\end{contrast}

\begin{contrast}{字节填充 / 零比特填充（每行依次为：适用场合·填充规则·典型协议）}
\citem{字节填充}{异步传输（逐个字符传送）}{把数据中的 0x7E 变为 0x7D 0x5E，把 0x7D 变为 0x7D 0x5D}{PPP（异步链路）}
\citem{零比特填充}{同步传输（连续比特流）}{每出现 5 个连续 1 就填入 1 个 0}{PPP（SONET/SDH 同步链路）、HDLC}
\end{contrast}

\begin{contrast}{链路 / 数据链路（每行依次为：组成·功能·示例）}
\citem{链路}{一段无源的点到点物理线路段}{仅提供物理传输通路}{一段光纤、一对双绞线}
\citem{数据链路}{链路 $+$ 实现通信协议的硬件和软件}{在物理链路上实现帧的传输与差错检测}{网络适配器（网卡）及其驱动程序}
\end{contrast}

\qtypename{一、单项选择题}{10 题，每题 3 分，共 30 分}
\q 数据链路层传输的数据单位是（\quad）
\choicesline{A. 比特}{B. 帧}{C. 分组}{D. 报文}

\q 数据链路层的三个基本问题是（\quad）
\choices{A. 封装成帧、透明传输、差错检测}{B. 路由选择、拥塞控制、流量控制}
{C. 调制、解调、复用}{D. 编号、确认、重传}

\q PPP 帧的标志字段值为（\quad）
\choicesline{A. 0x7E}{B. 0x7D}{C. 0xFF}{D. 0x03}

\q PPP 使用异步传输时，字节填充所用的转义符是（\quad）
\choicesline{A. 0x7E}{B. 0x7D}{C. 0x5E}{D. 0x03}

\q 零比特填充的规则是（\quad）
\choices{A. 每出现 5 个连续 1 后填入 1 个 0}{B. 每出现 5 个连续 0 后填入 1 个 1}
{C. 每出现 6 个连续 1 后填入 1 个 0}{D. 每出现 4 个连续 1 后填入 1 个 1}

\q CRC 检验采用（\quad）
\choicesline{A. 模 2 除法}{B. 普通十进制除法}{C. 加法求和}{D. 哈希运算}

\q 帧检验序列 FCS 的作用是（\quad）
\choices{A. 差错检测}{B. 差错纠正}
{C. 流量控制}{D. 路由选择}

\q 停止等待协议中，帧的序号只需 1 位，其原因是（\quad）
\choices{A. 只需要区分新帧和重传帧}{B. 为了提高发送速率}
{C. 1 位序号可以表示任意多个帧}{D. 为了节省帧的长度，与重传无关}

\q 回退 N 帧协议中，接收窗口的大小为（\quad）
\choicesline{A. 1}{B. 2}{C. 发送窗口的一半}{D. 与发送窗口相同}

\q PPP 协议的三个组成部分是（\quad）
\choices{A. 封装成帧、LCP、NCP}{B. 封装成帧、CRC、重传}
{C. IP、TCP、UDP}{D. 物理层、网络层、应用层}

\qtypename{二、多项选择题}{3 题，每题 4 分，共 12 分；每题有两个或两个以上正确项}
\q 数据链路层的三个基本问题包括（\quad）
\choicesline{A. 封装成帧}{B. 透明传输}{C. 差错检测}{D. 路由选择}

\q PPP 协议的组成部分包括（\quad）
\choicesline{A. 封装成帧}{B. 链路控制协议 LCP}{C. 网络控制协议 NCP}{D. 路由选择协议 RIP}

\q 关于 CRC 循环冗余检验，下列说法正确的有（\quad）
\choices{A. 收发双方必须使用相同的生成多项式}{B. 计算过程使用模 2 除法}
{C. 可以检测出多位比特差错}{D. 可以自动纠正所有差错}

\qtypename{三、判断题}{6 题，每题 2 分，共 12 分；正确的画“√”，错误的画“×”并写出改正后的表述}
\q 链路就是数据链路，二者没有区别。\quad（\quad）\quad 改正：\fixblank

\q 透明传输是指无论数据部分出现什么比特组合，都能按原样正确传输，不会与帧定界符混淆。\quad（\quad）\quad 改正：\fixblank

\q CRC 循环冗余检验既能检错，又能自动纠正所有差错。\quad（\quad）\quad 改正：\fixblank

\q 停止等待协议中，发送方每发送一个帧就等待确认，超时未收到确认则重传该帧。\quad（\quad）\quad 改正：\fixblank

\q 回退 N 帧协议的接收窗口大于 1。\quad（\quad）\quad 改正：\fixblank

\q PPP 异步传输用字节填充，同步传输用零比特填充。\quad（\quad）\quad 改正：\fixblank

\qtypename{四、计算题}{2 题，每题 5 分，共 10 分；要求写出计算步骤}
\q 要发送的数据为 101110，采用 CRC 的生成多项式 $G(x)=x^3+x+1$（即除数 1011）。求：（1）应添加的帧检验序列 FCS；（2）实际发送的比特序列。

\q 某数据链路层协议采用滑动窗口机制，帧序号字段为 3 位。分别求：（1）采用停止等待协议时的发送窗口大小；（2）采用连续 ARQ（回退 N 帧）时的最大发送窗口和接收窗口大小；（3）采用选择重传 ARQ 时的最大发送窗口和接收窗口大小。

\qtypename{五、简答题}{2 题，每题 8 分，共 16 分}
\q 简述数据链路层的三个基本问题，并说明透明传输的两种实现方法及其适用场合。\anslines{2.4cm}

\q 简述 PPP 协议的三个组成部分、PPP 帧的主要字段，以及 PPP 如何实现透明传输。\anslines{2.4cm}

\newpage
\qtypename{六、综合题}{1 题，共 20 分}
\q 两个相邻节点 A、B 之间采用数据链路层协议通信：链路数据率为 10 Mb/s，单向传播时延为 10 ms，帧长为 1000 bit，确认帧长为 100 bit，忽略处理时延和排队时延。请回答：（1）求数据帧的发送时延、单向传播时延和确认帧的发送时延；（2）若采用停止等待协议，求发送一帧并收到确认所需的时间以及信道利用率（列出公式）；（3）若采用连续 ARQ（回退 N 帧），发送窗口为 4，忽略确认帧发送时延，求信道利用率；（4）结合计算结果说明提高信道利用率的途径，并解释为什么在长时延链路上停止等待协议效率很低。\anslines{5cm}

\end{document}
